\section{QED and the Role of Masslessness for Renormalizability}

Quantum electrodynamics (QED) is one of the most successful theories in 
physics, providing predictions verified to extraordinary precision. A key 
reason for its consistency is the masslessness of the photon. This is not 
merely an empirical assumption but a structural requirement of the theory. 
Introducing a photon mass—even an extremely small one—breaks gauge symmetry 
and undermines the renormalizability and internal consistency of QED.

\subsection*{Gauge symmetry as the foundation of QED}

QED is built on the local $U(1)$ gauge symmetry of the electromagnetic 
field. This symmetry determines:

\begin{itemize}
	\item the form of the interaction between photons and charged particles,
	\item the structure of the electromagnetic current,
	\item the conservation of electric charge,
	\item and the cancellation of ultraviolet divergences.
\end{itemize}

The Ward–Takahashi identities, which are consequences of gauge symmetry, 
ensure that propagators, vertex functions and loop corrections maintain the 
proper relations required for a finite, renormalizable theory.

A photon mass term breaks gauge symmetry and invalidates these identities.

\subsection*{Failure of the Ward–Takahashi identities}

In QED with a massless photon, the Ward–Takahashi identities enforce 
relations of the form
\[
q_\mu \Gamma^\mu = S^{-1}(p+q) - S^{-1}(p),
\]
where $\Gamma^\mu$ is the vertex function and $S(p)$ the fermion propagator.
These relations guarantee the cancellation of divergences and protect the 
structure of the electric charge as a conserved quantity.

A photon mass term violates the gauge symmetry from which these identities 
follow. As a result:

\begin{itemize}
	\item divergences no longer cancel systematically,
	\item the renormalization of the charge becomes ill-defined,
	\item and the predictive power of the theory is lost.
\end{itemize}

Thus, renormalizable QED fundamentally requires $m_\gamma = 0$.

\subsection*{Longitudinal modes and non-renormalizable interactions}

A massive photon introduces a longitudinal polarization mode. In scattering 
processes involving this mode, amplitudes typically grow with energy.  
In a gauge-invariant theory such growth is cancelled by delicate 
interference between diagrams, but without gauge symmetry these 
cancellations fail.

The result is a loss of perturbative unitarity and the appearance of 
non-renormalizable divergences. High-energy scattering amplitudes 
ultimately violate probability bounds unless new physics is added to 
restore consistency.

This demonstrates that masslessness is not optional—it is required to 
maintain the quantum consistency of the theory.

\subsection*{Relation to the electroweak theory}

In the Standard Model, the photon emerges after spontaneous symmetry 
breaking from the electroweak gauge group. Its masslessness is protected 
by the unbroken electromagnetic $U(1)$ symmetry. Unlike the $W^\pm$ and 
$Z$ bosons, which acquire mass via the Higgs mechanism, the photon remains 
exactly massless because no symmetry-breaking term couples to the 
electromagnetic gauge field.

Any photon mass term would require a mechanism beyond the Standard Model, 
one that necessarily breaks the electromagnetic gauge symmetry and destroys 
the renormalizable structure of QED.

\subsection*{Conceptual implications}

The masslessness of the photon is therefore deeply embedded in the quantum 
structure of electromagnetism:

\begin{itemize}
	\item it guarantees renormalizability,
	\item it ensures charge conservation,
	\item it stabilizes the high-energy behaviour of scattering amplitudes,
	\item and it preserves the delicate cancellations required by gauge 
	symmetry.
\end{itemize}

These features are not optional; they are essential for the internal 
consistency of the theory. A massive photon would require either a new 
symmetry-breaking mechanism or a fundamentally different quantum field 
theory of electromagnetism.

\subsection*{Transition to the role of helicity and entanglement}

The next chapter examines an experimentally accessible fingerprint of 
photon masslessness: the helicity structure of quantum entanglement. 
The fact that photons exhibit only two helicity states is both a theoretical 
and experimental signature of a massless spin-1 particle. A massive photon 
would possess a third, longitudinal state, which would alter the structure 
of entangled-photon experiments in measurable ways.

Thus, the theoretical arguments from QED connect directly to observable 
quantum behaviour.
